\begin{figure}[t]
\centering
\begin{tikzpicture}[font=\scriptsize,x=1cm,y=1cm]
% Untiled
\node[anchor=west,font=\footnotesize\bfseries] at (0,3.45) {Untiled: one array per (slot, word) spanning all 385{,}024 workers};
\foreach \i/\t in {0/{slot 0, word 0},1/{slot 0, word 1},2/{\dots},3/{slot 15, word 4}}{
  \fill[cA!\the\numexpr 20+\i*15\relax] ({\i*3.4},2.6) rectangle ({\i*3.4+3.3},3.1);
  \node at ({\i*3.4+1.65},2.85) {\t};
}
\node[anchor=west,text=inkS] at (0,2.3) {a block's 256 threads touch 80 regions, each 1.5\,MB apart};
% Tiled
\node[anchor=west,font=\footnotesize\bfseries] at (0,1.65) {Tiled: each block's 256 workers keep all their slots and words together};
\foreach \b in {0,1,2}{
  \fill[cC!45] ({\b*4.5},0.8) rectangle ({\b*4.5+4.4},1.3);
  \node at ({\b*4.5+2.2},1.05) {block \b: 16 slots, 256 workers};
}
\node at (13.75,1.05) {\dots};
% Compact record
\node[anchor=west,font=\footnotesize\bfseries] at (0,0.3) {Compact element record for one (tile, slot): 4{,}352 bytes};
\fill[cC!70] (0,-0.6) rectangle (11.0,-0.1);
\node[text=white] at (5.5,-0.35) {256 low records of 16 bytes (4{,}096 B): 128 low bits of each worker's element};
\fill[cD] (11.05,-0.6) rectangle (13.85,-0.1);
\node at (12.45,-0.35) {256-byte tail plane};
\node[anchor=north west,text=inkS,align=left,text width=13.6cm] at (0,-0.75)
  {17 bytes per element instead of 20: the top 3 bits (or 6, with the cached branch) share one byte;
   the four state arrays fall from 470 to 399.5\,MiB.};
\end{tikzpicture}
\caption{How walk state is laid out in global memory. Tiling puts all of
  a block's state in one contiguous region while each warp still reads
  adjacent words ($+2.1\%$); the compact record then packs each 131-bit
  element into 17 bytes ($+6.3\%$). Neither changes a single bit of the
  walk.}
\label{fig:tiling}
\end{figure}
